\documentclass[11pt,a4paper]{article}

% Packages
\usepackage[utf8]{inputenc}
\usepackage[T1]{fontenc}
\usepackage[margin=2.5cm]{geometry}
\usepackage{graphicx}
\usepackage{xcolor}
\usepackage{listings}
\usepackage{hyperref}
\usepackage{fancyhdr}
\usepackage{tcolorbox}
\usepackage{enumitem}

% Code listing style
\lstset{
    language=C,
    basicstyle=\ttfamily\small,
    keywordstyle=\color{blue}\bfseries,
    commentstyle=\color{green!60!black}\itshape,
    stringstyle=\color{red},
    showstringspaces=false,
    breaklines=true,
    frame=single,
    numbers=left,
    numberstyle=\tiny\color{gray},
    tabsize=2
}

% Color boxes
\tcbuselibrary{skins,breakable}
\newtcolorbox{objectivebox}{colback=blue!5!white,colframe=blue!75!black,title=Lab Objectives,fonttitle=\bfseries}
\newtcolorbox{prereqbox}{colback=yellow!5!white,colframe=yellow!75!black,title=Prerequisites,fonttitle=\bfseries}
\newtcolorbox{solutionbox}{colback=green!5!white,colframe=green!75!black,title=Solution,fonttitle=\bfseries,breakable}

% Header and footer
\pagestyle{fancy}
\fancyhf{}
\lhead{USP Training Course 1}
\rhead{Lab 1.2: Send and Receive Messages}
\cfoot{\thepage}

\title{\textbf{Lab 1.2: Send and Receive Messages}\\
\large Course 1: LoRa \& LoRaWAN Fundamentals}
\author{USP Training Series}
\date{}

\begin{document}
\maketitle
\thispagestyle{fancy}

\section{Lab Overview}

\begin{objectivebox}
By the end of this lab, you will be able to:
\begin{itemize}
    \item Send unconfirmed uplink messages
    \item Format and encode payloads
    \item Receive and decode downlink messages
    \item Understand FPort usage
    \item Monitor message delivery in network server
\end{itemize}
\end{objectivebox}

\begin{prereqbox}
\textbf{Required:} Completion of Lab 1.1 (Basic LoRaWAN Join)\\
\textbf{Hardware:} Xiao nRF54L15 + LR1120 board
\end{prereqbox}

\section{Lab Duration}
\textbf{Estimated Time:} 45 minutes

\section{Background}

\subsection{LoRaWAN Ports}
FPort (Frame Port) identifies the application endpoint:
\begin{itemize}
    \item Port 0: Reserved for MAC commands
    \item Ports 1-223: Application data
    \item Port 224-255: Reserved for future use
\end{itemize}

\subsection{Payload Encoding}
Best practices for payload encoding:
\begin{itemize}
    \item Keep payloads small (< 51 bytes for SF12, < 222 bytes for SF7)
    \item Use binary encoding instead of ASCII
    \item Consider byte order (big-endian vs little-endian)
    \item Use The Things Network payload formatters for decoding
\end{itemize}

\section{Lab Instructions}

\subsection{Step 1: Implement Uplink Function}
Create a function to send sensor data as unconfirmed uplink.

\subsection{Step 2: Encode Sensor Data}
Pack temperature, humidity, and battery voltage into compact payload.

\subsection{Step 3: Send Periodic Uplinks}
Use Zephyr timers to send data every 60 seconds.

\subsection{Step 4: Receive Downlinks}
Implement callback to handle received downlink messages.

\subsection{Step 5: Test Communication}
Build, flash, and verify bidirectional communication.

\section{Solution}

\begin{solutionbox}
\subsection{Complete Implementation}

\begin{lstlisting}
#include <zephyr/kernel.h>
#include <zephyr/logging/log.h>
#include <zephyr/lorawan/lorawan.h>
#include <zephyr/random/random.h>

LOG_MODULE_REGISTER(lab_1_2, LOG_LEVEL_INF);

/* Simulated sensor readings */
static int16_t get_temperature(void)
{
    /* Simulate temperature: 20.5°C = 205 (x10) */
    return 205 + (sys_rand32_get() % 100) - 50;
}

static uint8_t get_humidity(void)
{
    /* Simulate humidity: 60% */
    return 60 + (sys_rand32_get() % 40);
}

static uint16_t get_battery_voltage(void)
{
    /* Simulate battery: 3.6V = 3600mV */
    return 3600 + (sys_rand32_get() % 200);
}

/* Encode sensor data into compact payload */
static int encode_sensor_data(uint8_t *payload)
{
    int16_t temp = get_temperature();
    uint8_t hum = get_humidity();
    uint16_t vbat = get_battery_voltage();

    /* Payload format (6 bytes):
     * [0-1]: Temperature (int16, x10, big-endian)
     * [2]:   Humidity (uint8, %)
     * [3-4]: Battery voltage (uint16, mV, big-endian)
     * [5]:   Sequence counter
     */
    static uint8_t seq = 0;

    payload[0] = (temp >> 8) & 0xFF;
    payload[1] = temp & 0xFF;
    payload[2] = hum;
    payload[3] = (vbat >> 8) & 0xFF;
    payload[4] = vbat & 0xFF;
    payload[5] = seq++;

    LOG_INF("Encoded: Temp=%.1f°C, Hum=%d%%, Vbat=%dmV",
            temp / 10.0, hum, vbat);

    return 6;
}

/* Send uplink message */
static void send_uplink(void)
{
    uint8_t payload[6];
    int len;
    int ret;

    len = encode_sensor_data(payload);

    ret = lorawan_send(2, payload, len,
                       LORAWAN_MSG_UNCONFIRMED);
    if (ret < 0) {
        LOG_ERR("Send failed: %d", ret);
    } else {
        LOG_INF("Uplink sent (%d bytes)", len);
        LOG_HEXDUMP_INF(payload, len, "Payload");
    }
}

/* Downlink callback */
static void dl_callback(uint8_t port, bool data_pending,
                       int16_t rssi, int8_t snr,
                       uint8_t len, const uint8_t *data)
{
    LOG_INF("========================================");
    LOG_INF("DOWNLINK RECEIVED");
    LOG_INF("========================================");
    LOG_INF("Port: %d", port);
    LOG_INF("RSSI: %d dBm", rssi);
    LOG_INF("SNR: %d dB", snr);
    LOG_INF("Pending: %s", data_pending ? "Yes" : "No");
    LOG_INF("Length: %d bytes", len);

    if (len > 0) {
        LOG_HEXDUMP_INF(data, len, "Data");

        /* Example: Parse simple command */
        if (port == 2 && len >= 1) {
            switch (data[0]) {
            case 0x01:
                LOG_INF("Command: LED ON");
                break;
            case 0x00:
                LOG_INF("Command: LED OFF");
                break;
            case 0xFF:
                LOG_INF("Command: Reset");
                break;
            default:
                LOG_INF("Command: Unknown (0x%02X)", data[0]);
                break;
            }
        }
    }
}

/* Uplink timer */
static void uplink_timer_handler(struct k_timer *timer)
{
    send_uplink();
}

K_TIMER_DEFINE(uplink_timer, uplink_timer_handler, NULL);

int main(void)
{
    int ret;
    const struct device *lora_dev;

    LOG_INF("==============================================");
    LOG_INF("Lab 1.2: Send and Receive Messages");
    LOG_INF("==============================================");

    lora_dev = DEVICE_DT_GET(DT_ALIAS(lora0));
    if (!device_is_ready(lora_dev)) {
        LOG_ERR("LoRa device not ready");
        return -1;
    }

    ret = lorawan_start();
    if (ret < 0) {
        LOG_ERR("Failed to start LoRaWAN: %d", ret);
        return ret;
    }

    /* Register downlink callback */
    struct lorawan_downlink_cb downlink_cb = {
        .port = LW_RECV_PORT_ANY,
        .cb = dl_callback
    };
    lorawan_register_downlink_callback(&downlink_cb);

    /* Join network */
    struct lorawan_join_config join_cfg = {
        .mode = LORAWAN_ACT_OTAA,
        .dev_eui = (uint8_t[]){0x00,0x00,0x00,0x00,
                               0x00,0x00,0x00,0x01},
        .otaa = {
            .join_eui = (uint8_t[]){0x00,0x00,0x00,0x00,
                                     0x00,0x00,0x00,0x00},
            .app_key = (uint8_t[]){0x00,0x00,0x00,0x00,
                                   0x00,0x00,0x00,0x00,
                                   0x00,0x00,0x00,0x00,
                                   0x00,0x00,0x00,0x00},
            .nwk_key = (uint8_t[]){0x00,0x00,0x00,0x00,
                                   0x00,0x00,0x00,0x00,
                                   0x00,0x00,0x00,0x00,
                                   0x00,0x00,0x00,0x00}
        }
    };

    LOG_INF("Starting join procedure...");
    ret = lorawan_join(&join_cfg);
    if (ret < 0) {
        LOG_ERR("Join failed: %d", ret);
        return ret;
    }

    /* Wait for join */
    while (!lorawan_is_joined()) {
        k_sleep(K_SECONDS(1));
    }

    LOG_INF("Network joined successfully!");

    /* Send first uplink immediately */
    send_uplink();

    /* Start periodic uplinks (60 seconds) */
    k_timer_start(&uplink_timer, K_SECONDS(60), K_SECONDS(60));

    LOG_INF("Periodic uplinks started (60s interval)");

    /* Main loop */
    while (1) {
        k_sleep(K_SECONDS(10));
    }

    return 0;
}
\end{lstlisting}

\subsection{TTN Payload Decoder}

Add this decoder function in your TTN application:

\begin{lstlisting}[language=JavaScript]
function decodeUplink(input) {
    var data = {};

    // Decode temperature (bytes 0-1, int16, x10)
    var temp_raw = (input.bytes[0] << 8) | input.bytes[1];
    if (temp_raw > 32767) temp_raw -= 65536;  // Handle negative
    data.temperature = temp_raw / 10.0;

    // Decode humidity (byte 2, uint8)
    data.humidity = input.bytes[2];

    // Decode battery voltage (bytes 3-4, uint16, mV)
    data.battery_voltage = ((input.bytes[3] << 8) |
                            input.bytes[4]) / 1000.0;

    // Sequence counter (byte 5)
    data.sequence = input.bytes[5];

    return {
        data: data,
        warnings: [],
        errors: []
    };
}
\end{lstlisting}

\end{solutionbox}

\section{Expected Output}

\begin{lstlisting}
[00:00:10.570] <inf> lab_1_2: Network joined successfully!
[00:00:10.571] <inf> lab_1_2: Encoded: Temp=20.7°C, Hum=65%, Vbat=3625mV
[00:00:10.572] <inf> lab_1_2: Uplink sent (6 bytes)
[00:00:10.573] <inf> lab_1_2: Payload: 00 cf 41 0e 29 00
[00:01:10.571] <inf> lab_1_2: Encoded: Temp=21.2°C, Hum=62%, Vbat=3638mV
[00:01:10.572] <inf> lab_1_2: Uplink sent (6 bytes)
[00:01:11.234] <inf> lab_1_2: ========================================
[00:01:11.235] <inf> lab_1_2: DOWNLINK RECEIVED
[00:01:11.236] <inf> lab_1_2: ========================================
[00:01:11.237] <inf> lab_1_2: Port: 2
[00:01:11.238] <inf> lab_1_2: RSSI: -52 dBm
[00:01:11.239] <inf> lab_1_2: SNR: 9 dB
[00:01:11.240] <inf> lab_1_2: Length: 1 bytes
[00:01:11.241] <inf> lab_1_2: Data: 01
[00:01:11.242] <inf> lab_1_2: Command: LED ON
\end{lstlisting}

\section{Deliverables}

\begin{enumerate}[itemsep=10pt]
    \item Source code with uplink/downlink implementation
    \item Serial log showing bidirectional communication
    \item TTN console screenshot with decoded payload
    \item Test downlink from TTN showing device receives it
\end{enumerate}

\section{Additional Challenges}

\subsection{Challenge 1: Payload Compression}
Reduce payload to 4 bytes using:
\begin{itemize}
    \item Temperature: 1 byte (range -40 to +85°C)
    \item Humidity: 1 byte (0-100%)
    \item Battery: 1 byte (2.0V to 4.2V)
    \item Counter: 1 byte
\end{itemize}

\subsection{Challenge 2: Downlink Actuator}
Control actual LED based on downlink commands.

\vfill
\noindent\rule{\textwidth}{0.4pt}
\begin{center}
\textit{USP Training Course 1 - Lab 1.2}\\
\textit{Version 1.0 - 2025}
\end{center}

\end{document}
